%=============================================================================
% MÓDULO 06: RESULTADOS (Expansão Final)
%=============================================================================
% Frontiers in Public Health — Resultados objetivos, sem interpretação
%   6.1 Desempenho dos modelos
%   6.2 Análise de justiça
%   6.3 Análise interseccional
%   6.4 Estabilidade da validação cruzada
%=============================================================================

\section{Resultados}

\subsection{Desempenho dos Modelos}

A Tabela~\ref{tab:performance} apresenta o desempenho discriminativo de 
todos os quatro modelos no conjunto de teste separado.

\begin{table}[H]
\centering
\caption{Comparação de Desempenho dos Modelos de Aprendizado de Máquina}
\begin{tabular}{lcccc}
\toprule
\textbf{Modelo} & \textbf{F1 VC (média±dp)} & \textbf{F1 Teste} & 
\textbf{AUC-ROC} & \textbf{Acurácia} \\
\midrule
Regressão Logística & 0.7929±0.0313 & 0.8272 & 0.8877 & 0.8052 \\
Random Forest & 0.7679±0.0243 & 0.7875 & 0.8731 & 0.7727 \\
Gradient Boosting & 0.7769±0.0221 & 0.7805 & 0.8543 & 0.7792 \\
\textbf{SVM} & \textbf{0.7859±0.0229} & \textbf{0.8313} & \textbf{0.8794} & \textbf{0.8117} \\
\bottomrule
\end{tabular}
\label{tab:performance}
\end{table}

O modelo SVM alcançou o maior F1-Score (0.8313) e AUC-ROC competitivo 
(0.8794), com desempenho estável na validação cruzada (VC: 0.7859$\pm$0.0229). 
A Regressão Logística demonstrou o melhor AUC-ROC (0.8877) mas menor 
F1-Score (0.8272), sugerindo um trade-off entre discriminação e desempenho 
específico da classe. Random Forest e Gradient Boosting mostraram desempenho 
comparável mas ligeiramente inferior em todas as métricas.

Os resultados da validação cruzada indicam que todos os modelos apresentaram 
estabilidade razoável, com desvios-padrão abaixo de 0.04 entre as dobras. 
Isso sugere que as diferenças de desempenho observadas não são atribuíveis 
a variação aleatória na divisão de dados, mas refletem diferenças genuínas 
nas capacidades dos modelos e sua interação com os dados.

\subsection{Análise de Justiça}

\subsubsection{Justiça Baseada em Etnia}

A Tabela~\ref{tab:ethnicity_fairness} apresenta métricas de justiça 
estratificadas por etnia para todos os modelos.

\begin{table}[H]
\centering
\caption{Métricas de Justiça por Etnia}
\begin{tabular}{lcccccc}
\toprule
\textbf{Modelo} & \textbf{EOD} & \textbf{DPD} & 
\textbf{TPR$_{\text{maj}}$} & \textbf{TPR$_{\text{min}}$} & 
\textbf{FPR$_{\text{maj}}$} & \textbf{FPR$_{\text{min}}$} \\
\midrule
Reg.\ Logística & 0.0687 & 0.0479 & 0.833 & 0.815 & 0.182 & 0.375 \\
Random Forest & 0.0306 & 0.0324 & 0.833 & 0.778 & 0.200 & 0.417 \\
Grad.\ Boosting & 0.0636 & 0.0201 & 0.812 & 0.704 & 0.164 & 0.333 \\
\textbf{SVM} & \textbf{0.0492} & \textbf{0.0148} & 0.812 & 0.741 & 0.182 & 0.375 \\
\bottomrule
\end{tabular}
\label{tab:ethnicity_fairness}
\end{table}

Todos os modelos apresentaram taxas de falso positivo mais altas para 
grupos minoritários (FPR: 0.333--0.417) em comparação com grupos 
majoritários (FPR: 0.164--0.200). O modelo Random Forest mostrou o 
menor EOD (0.0306) para etnia, enquanto o modelo SVM alcançou o menor 
DPD (0.0148), indicando que diferentes modelos apresentam diferentes 
perfis de justiça dependendo da métrica utilizada. Essa descoberta 
destaca a importância de avaliar múltiplas métricas de justiça em vez 
de depender de um único critério.

As taxas consistentemente mais altas de falso positivo para grupos 
minoritários em todos os modelos sugerem um viés sistemático na forma 
como pacientes minoritários são classificados. Esse padrão indica que 
os modelos são mais propensos a identificar incorretamente pacientes 
minoritários como tendo diabetes quando não têm, potencialmente levando 
a intervenções desnecessárias e custos de saúde.

\subsubsection{Justiça Baseada em Gênero}

A Tabela~\ref{tab:gender_fairness} apresenta métricas de justiça 
estratificadas por gênero.

\begin{table}[H]
\centering
\caption{Métricas de Justiça por Gênero}
\begin{tabular}{lcccccc}
\toprule
\textbf{Modelo} & \textbf{EOD} & \textbf{DPD} & 
\textbf{TPR$_{\text{M}}$} & \textbf{TPR$_{\text{F}}$} & 
\textbf{FPR$_{\text{M}}$} & \textbf{FPR$_{\text{F}}$} \\
\midrule
Reg.\ Logística & 0.0838 & 0.1010 & 0.833 & 0.815 & 0.182 & 0.375 \\
Random Forest & 0.1143 & 0.1010 & 0.833 & 0.778 & 0.200 & 0.417 \\
Grad.\ Boosting & 0.0787 & 0.0843 & 0.812 & 0.704 & 0.164 & 0.333 \\
\textbf{SVM} & \textbf{0.0566} & \textbf{0.0672} & 0.812 & 0.741 & 0.182 & 0.375 \\
\bottomrule
\end{tabular}
\label{tab:gender_fairness}
\end{table}

A análise baseada em gênero revelou padrões semelhantes, com pacientes 
femininas experimentando taxas de falso positivo mais altas em todos os 
modelos. O modelo SVM alcançou o menor EOD (0.0566) e DPD (0.0672) para 
gênero, indicando o desempenho mais equitativo nesta dimensão.

As disparidades baseadas em gênero observadas neste estudo são consistentes 
com pesquisas anteriores que mostram que modelos de predição clínica 
frequentemente apresentam desempenho diferente entre grupos de gênero. 
Essas diferenças podem refletir variações biológicas na apresentação da 
doença, diferenças nos padrões de utilização da saúde ou vieses na forma 
como dados clínicos são coletados e registrados para diferentes gêneros.

\subsection{Análise Interseccional}

A Tabela~\ref{tab:intersectional} apresenta métricas de desempenho na 
intersecção de gênero e etnia.

\begin{table}[H]
\centering
\caption{Análise Interseccional: Desempenho por Gênero $\times$ Etnia}
\begin{tabular}{lcccc}
\toprule
\textbf{Grupo} & \textbf{Acurácia} & \textbf{F1-Score} & 
\textbf{TPR} & \textbf{$n$} \\
\midrule
Majoritário--Feminino & 0.8000 & 0.8235 & 0.833 & 65 \\
Majoritário--Masculino & 0.8125 & 0.8276 & 0.833 & 38 \\
Minoritário--Feminino & 0.7826 & 0.7500 & 0.714 & 33 \\
Minoritário--Masculino & 0.8000 & 0.7273 & 0.750 & 18 \\
\bottomrule
\end{tabular}
\label{tab:intersectional}
\end{table}

A análise interseccional revelou que pacientes minoritários--femininos 
experimentaram o menor TPR (0.714) e F1-Score (0.750), indicando viés 
amplificado na intersecção de múltiplas identidades marginalizadas. 
Esse grupo também mostrou a maior lacuna de desempenho em relação ao 
grupo de melhor desempenho (majoritário--masculino), com diferença de 
TPR de 0.119 e diferença de F1-Score de 0.078.

Os achados interseccionais são particularmente preocupantes porque 
sugerem que mulheres minoritárias enfrentam desvantagem acumulada que 
não pode ser capturada analisando raça ou gênero isoladamente. Esse 
grupo apresentou tanto sensibilidade mais baixa (menos verdadeiros 
positivos) quanto acurácia mais baixa, indicando que o modelo é menos 
eficaz na identificação correta de diabetes nessa população. O menor 
tamanho da amostra para homens minoritários (n=18) pode ter contribuído 
para as diferenças de desempenho observadas, mas o padrão consistente 
em múltiplas métricas sugere uma preocupação genuína de justiça.

\subsection{Estabilidade da Validação Cruzada}

A Tabela~\ref{tab:cross_validation} resume os resultados da validação 
cruzada para todos os modelos.

\begin{table}[H]
\centering
\caption{Avaliação de Estabilidade da Validação Cruzada}
\begin{tabular}{lccc}
\toprule
\textbf{Modelo} & \textbf{F1 Médio} & \textbf{DP F1} & \textbf{Dobras VC} \\
\midrule
Regressão Logística & 0.7929 & 0.0313 & 5 \\
Random Forest & 0.7679 & 0.0243 & 5 \\
Gradient Boosting & 0.7769 & 0.0221 & 5 \\
SVM & 0.7859 & 0.0229 & 5 \\
\bottomrule
\end{tabular}
\label{tab:cross_validation}
\end{table}

Todos os modelos demonstraram desempenho estável nas dobras de validação 
cruzada, com desvios-padrão abaixo de 0.04. O modelo Gradient Boosting 
mostrou a maior estabilidade (DP: 0.0221), seguido de perto pelo SVM 
(DP: 0.0229). Esses resultados indicam que as métricas de desempenho e 
justiça observadas são confiáveis e não são atribuíveis a variação 
aleatória na divisão de dados.

Os resultados de estabilidade da validação cruzada são importantes para 
estabelecer confiança nos achados de justiça. Se os modelos mostrassem 
alta variância entre as doras, seria difícil determinar se as disparidades 
observadas eram genuínas ou artefatos de divisões específicas de dados. 
O desempenho consistente entre as doras fornece garantia de que as 
métricas de justiça refletem padrões estáveis nos dados.
